% Generated draft: requires mathematical and source review.
\documentclass[11pt]{article}
\usepackage[margin=25mm]{geometry}
\usepackage{fontspec}
\setmainfont{Latin Modern Roman}
\usepackage{amsmath,amssymb,mathtools,unicode-math}
\setmathfont{Latin Modern Math}
\usepackage{enumitem,needspace,xurl,hyperref}
\setlength{\parindent}{0pt}
\setlength{\parskip}{5pt}
\setlength{\emergencystretch}{3em}
\newcommand{\N}{\mathbb{N}}
\newcommand{\Z}{\mathbb{Z}}
\newcommand{\Q}{\mathbb{Q}}
\newcommand{\R}{\mathbb{R}}
\newcommand{\C}{\mathbb{C}}
\newcommand{\F}{\mathbb{F}}
\newcommand{\A}{\mathbb{A}}
\newcommand{\PP}{\mathbb{P}}
\newcommand{\E}{\mathbb{E}}
\newcommand{\eps}{\varepsilon}
\newcommand{\dd}{\,\mathrm{d}}
\begin{document}
\section*{UnsolvedMath ID: 6200098}
\textbf{Source problem code:} catalogue entry
\subsection{Definitions and mathematical statement}
A word-hyperbolic group is a finitely generated group whose Cayley graph for a finite generating set is Gromov-hyperbolic. A virtually cyclic group has a cyclic subgroup of finite index, including the finite-group case. A nontrivial splitting over virtually cyclic groups can be expressed as an action on a Bass--Serre tree with no global fixed point and with virtually cyclic edge stabilizers.

Let $G$ be word-hyperbolic and admit no such splitting. Can there be subgroups $H,K\leq G$ such that
\[
[G:H]=\infty,\qquad [G:K]<\infty,\qquad H\cong K
\]
as abstract groups? Equivalently, is it possible under the nonsplitting hypothesis for an infinite-index subgroup to be abstractly isomorphic to a finite-index subgroup? Neither torsion-freeness nor quasiconvexity of $H$ is assumed, and the isomorphism need not be induced by conjugation or by an automorphism of $G$.
\subsection{Short English statement}
Can a hyperbolic group with no virtually-cyclic splitting have abstractly isomorphic subgroups of infinite and finite index?
\subsection{Sources}
\begin{itemize}
\item Misha Kapovich, \emph{Problems on Boundaries of Groups and Kleinian Groups} (2005), Problem 98, PDF page 24, attributed to Kevin Whyte, \url{https://www.math.ucdavis.edu/~kapovich/EPR/problems.pdf}; accessed 2026-07-29, PDF-text-normalized extraction.
\item Local registry \texttt{data/unsolvedmath/problems.json}, permanent UnsolvedMath ID $6200098$, source code catalogue entry, \emph{Boundaries of Groups and Kleinian Groups --- Problem 98}. Dataset attribution: Ulam AI and the UnsolvedMath Contributors.
\end{itemize}
\end{document}
